% section: 基本4パターン

\section{基本4パターン}
  まずは多くの漸化式の最終的な終着地点になる,基本4パターンを紹介する.
  多くの可求型漸化式はこの形に持っていくことでその一般項を得る.
  その中でも特性方程式型に帰着後,等比型に帰着させることで一般項を求めることが多い.
  他の基本パターンも後の応用パターンの礎となる発想をもつため,この節は特に丁寧に取り組むことを推奨する.
  
  そして,蛇足にはなるが,一般論として「基礎」「基本」などの文言は簡単を意味しない.
  理論の根幹を成すという意味である.
  軽んじてはならない.

\input{第3章 漸化式の解説/03_基本4パターン/01_等差型/等差型.tex}
\input{第3章 漸化式の解説/03_基本4パターン/02_等比型/等比型.tex}
% subsection: 階差型 \, $a_{n+1}=a_n+f(n)$

  \subsection{階差型 \, $a_{n+1}=a_n+f(n)$}
    階差型は先の2つに比べて一段階難易度が高い.
    しかしきちんと先程の議論を理解して,丁寧に式を追うことで理解することができるものである.
    兎にも角にも解法を見てみよう.
    \begin{align*}
      &a_{n+1}=a_n+f(n)\\
      &a_{n+1}=\{a_{n-1}+f(n-1)\}+f(n)\\
      &a_{n+1}=\left[\{a_{n-2}+f(n-2)\}+f(n-1)\right]+f(n)\\
      &\qquad \vdots\\
      &a_{n+1}=a_1+f(1)+f(2)+\cdots+f(n)\\
      &a_{n+1}=a_1+\sum_{k=1}^{n} f(k)\\
      &\therefore \quad a_n=a_1+\sum_{k=1}^{n-1} f(k)
    \end{align*}
    $f(n)$によってこの後の処理が変わるものの,漸化式の解法としてはここまでである.
    ここからはシグマの計算の能力が問われる話で,この後にやることは等差数列の総和を求めることであったり,等比数列の総和を求めることであったり,あるいは部分分数分解を用いて解くことなどが考えられる.

    \noindent \textbf{【例題】}
    \begin{enumerate}[label=(\arabic*),parsep=3pt]
      \item $\displaystyle a_1=1,\, a_{n+1}=a_n+n$
      \item $\displaystyle a_1=1,\, a_{n+1}=a_n+2^n$
      \item $\displaystyle a_1=0,\, a_{n+1}=a_n+\frac{1}{n(n+1)}$
      \item $\displaystyle a_1=1,\, a_{n+1}=a_n+(n+1)^2$
    \end{enumerate}

    \noindent\textbf{【方針】}
    \begin{enumerate}[label=(\arabic*),parsep=3pt]
      \item 最もシンプルな形の階差数列です.先程の議論を踏まえて解きましょう.
      \item もちろん,\,$f(n)$となる関数には指数関数も含まれます.何にせよやることは\,(1)と同じです.
      \item シグマは部分分数分解で処理しましょう.
      \item 勘の良い人であれば式を見た瞬間に答えがわかるかもしれませんね.これは"あの数列"の総和と同じ見た目になりませんか?.
    \end{enumerate}

    \noindent\textbf{【解答】}
    \begin{enumerate}[label=(\arabic*),parsep=3pt]
      \item $\displaystyle a_n=1+\frac{n(n-1)}{2}$
      \item $\displaystyle a_n=2^n-1$
      \item $\displaystyle a_n=1-\frac{1}{n}$
      \item $\displaystyle a_n=\frac{1}{6}n(n+1)(2n+1)$
    \end{enumerate}

    \noindent\textbf{【解法】}
    \begin{enumerate}[label=(\arabic*),parsep=3pt]
      \item $\displaystyle a_1=1,\,a_{n+1}=a_n+n$\\
            階差型の一般項
            \[
              a_n=a_1+\sum_{k=1}^{n-1}f(k)
            \]
            を用いると,
            \begin{align*}
              a_n
                &=1+\sum_{k=1}^{n-1}k\\
                &=1+\frac{n(n-1)}{2}
            \end{align*}

      \item $\displaystyle a_1=1,\,a_{n+1}=a_n+2^n$\\
            同様に,
            \begin{align*}
              a_n
                &=1+\sum_{k=1}^{n-1}2^k\\
                &=1+(2+2^2+\cdots+2^{n-1})\\
                &=1+(2^n-2)\\
                &=2^n-1
            \end{align*}

      \item $\displaystyle a_1=0,\,a_{n+1}=a_n+\frac{1}{n(n+1)}$\\
            部分分数分解を用いると,
            \[
              \frac{1}{k(k+1)}
              =\frac{1}{k}-\frac{1}{k+1}
            \]
            であるから,
            \begin{align*}
              a_n
                &=\sum_{k=1}^{n-1}\frac{1}{k(k+1)}\\
                &=\sum_{k=1}^{n-1}
                  \left(\frac{1}{k}-\frac{1}{k+1}\right)\\
                &=\left(1-\frac12\right)
                  +\left(\frac12-\frac13\right)\\
                &\quad +\cdots+\left(\frac{1}{n-1}-\frac1n\right)\\
                &=1-\frac1n
            \end{align*}

      \item $\displaystyle a_1=1,\,a_{n+1}=a_n+(n+1)^2$\\
            階差型の一般項を用いると,
            \begin{align*}
              a_n
                &=1+\sum_{k=1}^{n-1}(k+1)^2\\
                &=1+\sum_{j=2}^{n}j^2\\
                &=\sum_{j=1}^{n}j^2\\
                &=\frac{1}{6}n(n+1)(2n+1)
            \end{align*}
    \end{enumerate}

\input{第3章 漸化式の解説/03_基本4パターン/04_階比型/階比型.tex}
